% GENERATED FILE. Edit canonical lessons, then run npm run generate:books.
% canonical-source-hash: fnv1a64:7a2d564bab7f1452
% canonical-lessons: AR-C71-mirah, AR-C71-misbah, AR-C71-miftah, AR-C71-sahn, AR-C71-sikkin

\chapter{Mirror, Lamp, Key, Plate, Knife}
\label{ch:ar-mirror-lamp-key-plate-knife}
\hlchaptermodality{71}

\begin{chapteropening}
\textbf{By the end of this chapter:} \emph{I can say a mirror, a lamp, a key, a plate and a knife.}

Complete the last lesson of chapter 71: I can say a mirror, a lamp, a key, a plate and a knife.
\end{chapteropening}

\section[\texorpdfstring{\ar{مرآة} (mirʾāh)}{مرآة (mirʾāh)}]{\ar{مرآة} (mirʾāh) — a mirror}

\label{lesson:AR-C71-mirah}



\begin{quote}
Before the new one: say the Arabic for a bed, then the Arabic for the floor; the earth.
\end{quote}

\subsection*{You'll want to know: \ar{مرآة}}
\textbf{\ar{مرآة}} — \emph{mirʾāh} — \textquotedblleft{}a mirror\textquotedblright{}.

\begin{cousinweb}
\textbf{\ar{مرآة}} (\emph{mirʾāh}) is built on the native Arabic root \emph{r-ʾ-y}, \textquotedblleft{}to see\textquotedblright{} (Lane). \emph{Mirʾāh} is an instrument noun, \textquotedblleft{}a thing to see with\textquotedblright{}.
\end{cousinweb}

\begin{grammarlens}[title={What you've built}]
One more word for this chapter.
\end{grammarlens}

\subsection*{Your turn}
\begin{itemize}
\raggedright
  \item \emph{Say it:} \emph{mirʾāh}
  \item \emph{Say it:} \emph{mirʾāh}, once more
  \item \emph{Say it:} say \emph{mirʾāh}
  \item \emph{Recall:} say the Arabic for a bed, then the Arabic for the floor; the earth, then say \emph{mirʾāh} again
\end{itemize}

\begin{tcolorbox}[breakable,colback=teal!4,colframe=teal!35!black,title={Before you move on}]
What does \ar{مرآة} mean? (\textquotedblleft{}A mirror\textquotedblright{}.) Say it once more.
\end{tcolorbox}
\section[\texorpdfstring{\ar{مصباح} (miṣbāḥ)}{مصباح (miṣbāḥ)}]{\ar{مصباح} (miṣbāḥ) — a lamp}

\label{lesson:AR-C71-misbah}



\begin{quote}
Before the new one: say the Arabic for the floor; the earth, then the Arabic for a mirror.
\end{quote}

\subsection*{You'll want to know: \ar{مصباح}}
\textbf{\ar{مصباح}} — \emph{miṣbāḥ} — \textquotedblleft{}a lamp\textquotedblright{}.

\begin{cousinweb}
\textbf{\ar{مصباح}} (\emph{miṣbāḥ}) is built on the native Arabic root \emph{ṣ-b-ḥ}, \textquotedblleft{}morning\textquotedblright{} (Lane). \emph{Miṣbāḥ} is an instrument noun on the root of morning, a thing that gives light.
\end{cousinweb}

\begin{grammarlens}[title={What you've built}]
One more word for this chapter.
\end{grammarlens}

\subsection*{Your turn}
\begin{itemize}
\raggedright
  \item \emph{Say it:} \emph{miṣbāḥ}
  \item \emph{Say it:} \emph{miṣbāḥ}, once more
  \item \emph{Say it:} say \emph{miṣbāḥ}
  \item \emph{Recall:} say the Arabic for the floor; the earth, then the Arabic for a mirror, then say \emph{miṣbāḥ} again
\end{itemize}

\begin{tcolorbox}[breakable,colback=teal!4,colframe=teal!35!black,title={Before you move on}]
What does \ar{مصباح} mean? (\textquotedblleft{}A lamp\textquotedblright{}.) Say it once more.
\end{tcolorbox}
\section[\texorpdfstring{\ar{مفتاح} (miftāḥ)}{مفتاح (miftāḥ)}]{\ar{مفتاح} (miftāḥ) — a key}

\label{lesson:AR-C71-miftah}



\begin{quote}
Before the new one: say the Arabic for a mirror, then the Arabic for a lamp.
\end{quote}

\subsection*{You'll want to know: \ar{مفتاح}}
\textbf{\ar{مفتاح}} — \emph{miftāḥ} — \textquotedblleft{}a key\textquotedblright{}.

\begin{cousinweb}
\textbf{\ar{مفتاح}} (\emph{miftāḥ}) is built on the native Arabic root \emph{f-t-ḥ}, \textquotedblleft{}to open\textquotedblright{} (Lane). \emph{Miftāḥ} is the \emph{mifʿāl} instrument noun, \textquotedblleft{}an instrument of opening\textquotedblright{}.
\end{cousinweb}

\begin{grammarlens}[title={What you've built}]
One more word for this chapter.
\end{grammarlens}

\subsection*{Your turn}
\begin{itemize}
\raggedright
  \item \emph{Say it:} \emph{miftāḥ}
  \item \emph{Say it:} \emph{miftāḥ}, once more
  \item \emph{Say it:} say \emph{miftāḥ}
  \item \emph{Recall:} say the Arabic for a mirror, then the Arabic for a lamp, then say \emph{miftāḥ} again
\end{itemize}

\begin{tcolorbox}[breakable,colback=teal!4,colframe=teal!35!black,title={Before you move on}]
What does \ar{مفتاح} mean? (\textquotedblleft{}A key\textquotedblright{}.) Say it once more.
\end{tcolorbox}
\section[\texorpdfstring{\ar{صحن} (ṣaḥn)}{صحن (ṣaḥn)}]{\ar{صحن} (ṣaḥn) — a plate}

\label{lesson:AR-C71-sahn}



\begin{quote}
Before the new one: say the Arabic for a lamp, then the Arabic for a key.
\end{quote}

\subsection*{You'll want to know: \ar{صحن}}
\textbf{\ar{صحن}} — \emph{ṣaḥn} — \textquotedblleft{}a plate\textquotedblright{}.

\begin{cousinweb}
\textbf{\ar{صحن}} (\emph{ṣaḥn}) is built on the native Arabic root \emph{ṣ-ḥ-n}, \textquotedblleft{}a large bowl; an open court\textquotedblright{} (Lane). Lane gives a large bowl, and also the open court of a house.
\end{cousinweb}

\begin{grammarlens}[title={What you've built}]
One more word for this chapter.
\end{grammarlens}

\subsection*{Your turn}
\begin{itemize}
\raggedright
  \item \emph{Say it:} \emph{ṣaḥn}
  \item \emph{Say it:} \emph{ṣaḥn}, once more
  \item \emph{Say it:} say \emph{ṣaḥn}
  \item \emph{Recall:} say the Arabic for a lamp, then the Arabic for a key, then say \emph{ṣaḥn} again
\end{itemize}

\begin{tcolorbox}[breakable,colback=teal!4,colframe=teal!35!black,title={Before you move on}]
What does \ar{صحن} mean? (\textquotedblleft{}A plate\textquotedblright{}.) Say it once more.
\end{tcolorbox}
\section[\texorpdfstring{\ar{سكين} (sikkīn)}{سكين (sikkīn)}]{\ar{سكين} (sikkīn) — a knife}

\label{lesson:AR-C71-sikkin}



\begin{quote}
Before the new one: say the Arabic for a key, then the Arabic for a plate.
\end{quote}

\subsection*{You'll want to know: \ar{سكين}}
\textbf{\ar{سكين}} — \emph{sikkīn} — \textquotedblleft{}a knife\textquotedblright{}.

\begin{cousinweb}
\textbf{\ar{سكين}} (\emph{sikkīn}) is usually traced to Aramaic \emph{sakkīnā}, \textquotedblleft{}knife\textquotedblright{} (EtymDB), with Hebrew \emph{sakkin} as a cousin. Lane lists it under \emph{s-k-n}, \textquotedblleft{}to be still\textquotedblright{}, and quotes the old explanation that a knife stills the animal it slaughters; that is a guess of the Arabic grammarians, a folk etymology.
\end{cousinweb}

\begin{grammarlens}[title={What you've built}]
One more word for this chapter.
\end{grammarlens}

\subsection*{Your turn}
\begin{itemize}
\raggedright
  \item \emph{Say it:} \emph{sikkīn}
  \item \emph{Say it:} \emph{sikkīn}, once more
  \item \emph{Say it:} say \emph{sikkīn}
  \item \emph{Recall:} say the Arabic for a key, then the Arabic for a plate, then say \emph{sikkīn} again
\end{itemize}

\begin{tcolorbox}[breakable,colback=teal!4,colframe=teal!35!black,title={Before you move on}]
What does \ar{سكين} mean? (\textquotedblleft{}A knife\textquotedblright{}.) Say it once more.
\end{tcolorbox}
